% !TEX TS-program = xelatex
\documentclass[11pt,a4paper]{article}

\usepackage[a4paper,margin=2.2cm]{geometry}
\usepackage{fontspec}
\usepackage{xeCJK}
\usepackage{graphicx}
\usepackage{xcolor}
\usepackage{hyperref}
\usepackage{titlesec}
\usepackage{enumitem}
\usepackage{fancyhdr}
\usepackage{setspace}

\setmainfont{Times}
\setsansfont{Helvetica}
\setmonofont{Menlo}
\IfFontExistsTF{PingFang SC}{
  \setCJKmainfont{PingFang SC}
  \setCJKmonofont{PingFang SC}
}{
  \IfFontExistsTF{Songti SC}{
    \setCJKmainfont{Songti SC}
    \setCJKmonofont{Songti SC}
  }{
    \IfFontExistsTF{STSong}{
      \setCJKmainfont{STSong}
      \setCJKmonofont{STSong}
    }{
      \IfFontExistsTF{SimSun}{
        \setCJKmainfont{SimSun}
        \setCJKmonofont{SimSun}
      }{
        \setCJKmainfont{Arial Unicode MS}
        \setCJKmonofont{Arial Unicode MS}
      }
    }
  }
}

\definecolor{Accent}{HTML}{2E5E4E}
\definecolor{SoftAccent}{HTML}{8FB09E}
\definecolor{Muted}{HTML}{555555}

\hypersetup{
  colorlinks=true,
  linkcolor=Accent,
  urlcolor=Accent
}

\titleformat{\section}
  {\Large\bfseries\color{Accent}}{\thesection}{0.6em}{}
\titleformat{\subsection}
  {\bfseries\color{Accent}}{\thesubsection}{0.5em}{}
\titlespacing*{\section}{0pt}{1.2em}{0.5em}
\titlespacing*{\subsection}{0pt}{0.8em}{0.3em}

\setlist[itemize]{leftmargin=*,topsep=4pt,itemsep=2pt}
\setlist[enumerate]{leftmargin=*,topsep=4pt,itemsep=2pt}

\pagestyle{fancy}
\fancyhf{}
\lhead{BioPath}
\rhead{Technical Report}
\cfoot{\thepage}
\setlength{\headheight}{14pt}

\newcommand{\CN}{\textbf{中文：}}
\newcommand{\EN}{\textbf{English:}}

\begin{document}
\onehalfspacing

\begin{titlepage}
  \vspace*{1.2cm}
  \noindent{\Large \textbf{BioPath}}\par
  \vspace{0.4cm}
  {\huge\bfseries 技术报告 / Technical Report}\par
  \vspace{0.6cm}
  {\large Bio-inspired Optimization for Precision Pest Management}\par
  \vspace{0.8cm}
  \textcolor{SoftAccent}{\rule{\linewidth}{1.2pt}}\par
  \vspace{0.8cm}
  \textbf{Version:} 1.0\par
  \textbf{Date:} \today\par
  \vfill
  \noindent
  \textcolor{Muted}{This report is bilingual (Chinese/English) and describes the BioPath MVP
  implementation, algorithms, and Web UI usage.}
\end{titlepage}

\section{概览 / Overview}
\CN BioPath 是一个轻量级优化工具，用于在栅格地图上推荐最优布点位置，
以最小化到最近陷阱的平均行走距离，并支持基于风险权重的优化。
系统包含 CLI 与本地 Web 可视化界面，能够输出指标、热力图与报告。\par
\EN BioPath is a lightweight optimization tool that recommends trap locations on grid maps,
minimizing the average walking distance to the nearest trap and supporting risk-weighted
optimization. The MVP includes a CLI and a local Web UI, and outputs metrics, heatmaps,
and reports.

\section{架构 / Architecture}
\CN 系统采用模块化结构，核心组件包括：\par
\begin{itemize}
  \item 地图解析与数据模型：\texttt{biopath/mapio.py}
  \item 候选点生成：\texttt{biopath/candidates.py}
  \item 目标函数与指标：\texttt{biopath/objective.py}
  \item 贪心优化与局部改进：\texttt{biopath/optimizer.py}
  \item 可视化与报告：\texttt{biopath/viz.py}, \texttt{biopath/report.py}
  \item CLI：\texttt{biopath/cli.py}
  \item Web UI：\texttt{biopath/web.py} 与 \texttt{biopath/web/}
\end{itemize}
\EN The system is modular. Key components include map parsing, candidate generation,
objective and metrics, greedy optimization with local swaps, visualization/reporting,
CLI entry point, and a local Web UI server with static assets.

\section{数据模型 / Data Model}
\CN 地图采用 JSON 格式，包含：\par
\begin{itemize}
  \item \texttt{name}: 地图名称
  \item \texttt{cell\_size\_m}: 单元格尺寸（米）
  \item \texttt{ascii}: ASCII 栅格，\texttt{.} 为可行走，\texttt{\#} 为障碍
  \item \texttt{weights}（可选）：与栅格同尺寸的权重矩阵
\end{itemize}
\CN 若未提供 \texttt{weights}，所有可行走单元权重默认 1.0。\par
\EN The map JSON includes a name, cell size, ASCII grid, and optional weights matrix.
If weights are omitted, all walkable cells default to weight 1.0.

\section{算法与实现 / Algorithms}
\subsection{距离图 / Distance Map}
\CN 采用多源 BFS 在 4 邻域上计算每个可行走单元到最近陷阱的最短距离。
距离单位为米（按 \texttt{cell\_size\_m} 累加）。\par
\EN A multi-source BFS over the 4-neighborhood computes the shortest distance from every
walkable cell to its nearest trap, measured in meters.

\subsection{目标函数 / Objective}
\CN 提供两种目标：\par
\begin{itemize}
  \item \textbf{Mean}: 所有可行走单元的平均距离
  \item \textbf{Weighted Mean}: 以权重加权的平均距离
\end{itemize}
\EN Two objectives are supported: unweighted mean distance and weighted mean distance
using the optional weights matrix.

\subsection{贪心优化 / Greedy Optimization}
\CN 算法迭代选择陷阱位置，每次选择能最小化目标的候选点，支持局部交换改进。
适合中小规模地图的快速求解。\par
\EN The solver greedily selects traps, optionally applying single-swap local improvements.
It is fast and suitable for small-to-medium maps.

\section{指标 / Metrics}
\CN 输出的指标包括：平均距离、加权平均距离、最大距离、P95 距离，
以及可选的覆盖率（给定阈值半径）。\par
\EN Metrics include mean distance, weighted mean distance, max distance, P95 distance,
and optional coverage within a user-specified radius.

\section{Web UI 使用说明 / Web UI Usage}
\CN \textbf{启动步骤：}\par
\begin{enumerate}
  \item 在项目根目录运行：\texttt{python3 -m biopath.web --port 8000}
  \item 打开浏览器：\url{http://127.0.0.1:8000}
  \item 在 \textbf{Map Input} 选择样例或上传 JSON
  \item 在 \textbf{Map Builder} 创建空地图或绘制障碍/权重
  \item 设置参数并点击 \textbf{Run Optimization}
  \item 需要导出时使用 \textbf{Download map/results}
\end{enumerate}
\EN \textbf{Steps:}\par
\begin{enumerate}
  \item From the project root: \texttt{python3 -m biopath.web --port 8000}
  \item Open: \url{http://127.0.0.1:8000}
  \item Load a sample/upload JSON in \textbf{Map Input}
  \item Use \textbf{Map Builder} to create or paint obstacles/weights
  \item Configure settings and click \textbf{Run Optimization}
  \item Export via \textbf{Download map/results}
\end{enumerate}

\section{错误处理与验证 / Validation}
\CN 系统对地图格式、权重、参数范围进行校验；若存在不可达单元，
目标函数返回无穷大，UI 会以灰色显示。\par
\EN The system validates input maps, weights, and parameters. Unreachable cells yield
infinite objectives and are rendered in gray.

\section{限制与后续 / Limitations \& Next Steps}
\CN 目前采用贪心优化，尚非全局最优；后续可引入更强的搜索策略或多目标约束，
并支持 GIS 数据与更复杂场景。\par
\EN Greedy optimization is not globally optimal. Future work can include stronger search
methods, multi-objective constraints, and GIS-based inputs.

\end{document}
